\documentclass[11pt,a4paper]{article}

\usepackage[margin=1in]{geometry}
\usepackage{amsmath,amssymb}
\usepackage{booktabs}
\usepackage{array}
\usepackage{xcolor}
\usepackage{graphicx}
\usepackage{caption}
\usepackage[colorlinks=true,linkcolor=blue!50!black,urlcolor=blue!50!black,
            citecolor=blue!50!black]{hyperref}
\usepackage{tcolorbox}
\usepackage{enumitem}
\usepackage{fancyhdr}

\newcommand{\code}[1]{\texttt{\small #1}}
\newcommand{\fhat}{\hat{f}}

\pagestyle{fancy}
\fancyhf{}
\rhead{\small Online adaptive correction of a frozen VLA}
\lhead{\small Internal report}
\cfoot{\thepage}
\renewcommand{\headrulewidth}{0.4pt}

\title{\textbf{Online Adaptive Correction of a Frozen\\Vision--Language--Action Model}\\[0.4em]
\large Theory, method, and results}
\author{Taheri et al.}
\date{1 September 2026}

\begin{document}
\maketitle
\thispagestyle{fancy}

\begin{tcolorbox}[colback=blue!4,colframe=blue!45!black,title=\textbf{Result in one line}]
A fault at the Cartesian action interface of a frozen 3.35\,B vision--language--action model is
identified and corrected \emph{within a single episode}, from the robot's own
motion, with no gradients through the task and no search over task success ---
\textbf{across four LIBERO suites, 29\% $\rightarrow$ 73\% pooled,
$p = 1.1\times 10^{-7}$} --- \textbf{provided the correction is restricted to
the channels on which the fault is identifiable.} Correcting the rest is worse
than doing nothing.
\end{tcolorbox}

\begin{tcolorbox}[colback=gray!6,colframe=gray!55!black]
\textbf{The scope condition is the contribution.} With all six action channels
corrected, the method is significant on \emph{one suite of four}. Restricted to
identifiable channels it is significant on \emph{all four}. Identifiability is
predictable before deployment from healthy data alone.
\end{tcolorbox}

\vspace{0.4em}
\noindent\textbf{Artifacts.} Video: \code{results/phase05/adaptive\_vs\_frozen.mp4}.
Data: \code{results/}. Code: \code{openpi/\{error\_signal, openloop\_id,
adaptive\_law, compare\_video\}.py}. Full theory:
\code{docs/ADAPTIVE\_CONTROL\_VLA.md}.

\noindent\textbf{Status.} Simulation; hardware validation in progress.
Section~\ref{sec:limits} states what this does not establish.

\tableofcontents

%==========================================================================
\section{The problem}

A vision--language--action (VLA) policy deployed on a robot degrades when the
hardware does: an actuator develops a constant offset, a joint loses gain, a
sensor drifts. The policy was trained on healthy hardware and systematically
commands the wrong thing.

Retraining a 3.35\,B-parameter policy is not possible mid-deployment, and often
not possible at all --- the operator typically owns neither the weights nor the
training data. Black-box search over task success avoids retraining but discards
everything known about the plant and pays for it in episodes.

The observation this rests on: \emph{the robot measures its own end-effector
pose every control step.} If the fault is observable there it can be identified
directly, inside one episode, without consulting the task reward.

The finding that organises everything below is that ``observable there'' is
\textbf{channel-dependent}, decidable in advance, and that acting on channels
which fail the test is actively harmful.

%==========================================================================
\section{Setting}

\begin{itemize}[leftmargin=1.4em,itemsep=2pt]
  \item \textbf{Model.} $\pi_{0.5}$-LIBERO, frozen: 3.35\,B parameters, no
        fine-tuning, no backpropagation through the environment. Faults are
        injected client-side between the policy and \code{env.step}; the model is
        never modified.
  \item \textbf{Benchmarks.} \code{libero\_spatial} (nominal \textbf{99.0\%} over
        500 episodes), plus \code{libero\_goal}, \code{libero\_object},
        \code{libero\_10}.
  \item \textbf{Controller.} OSC\_POSE at 20\,Hz, \code{control\_delta = True},
        \code{output\_max} $= [0.05,0.05,0.05,0.5,0.5,0.5]$ (m, m, m, rad, rad, rad).
  \item \textbf{Controls.} Frozen and adaptive arms share task, initial state,
        policy and fault; the only difference is whether the law runs.
\end{itemize}

%==========================================================================
\section{Fault screening}
\label{sec:screen}

Nine conditions screened (20 episodes each) for one that hurts enough to matter
and leaves the robot recoverable.

\begin{center}
\begin{tabular}{llr>{\raggedright\arraybackslash}p{5.0cm}}
\toprule
\textbf{Fault} & \textbf{Severity} & \textbf{Success} & \textbf{Verdict} \\
\midrule
nominal      & ---   & 95\,\%  & reference \\
brightness   & 0.1 / 0.2 / 0.4 & 100\,\% & no effect at any level \\
gain         & 0.7   & 95\,\%  & no effect \\
gain         & 0.5   & 75\,\%  & mild \\
gain         & 0.3   & 0\,\%   & catastrophic \\
\textbf{offset} & \textbf{0.05} & \textbf{55\,\%} & \textbf{degraded, recoverable} \\
offset       & 0.10  & 5\,\%   & near-lethal \\
offset       & 0.20  & 0\,\%   & catastrophic \\
\bottomrule
\end{tabular}
\end{center}

Visual perturbation does nothing at any level tested; actuator faults show a
sharp cliff. Offset $0.05$ --- a constant $f = +0.05$ added to the six arm
dimensions of every action, i.e.\ a miscalibrated actuator --- is the primary
condition throughout.

%==========================================================================
\section{Is the fault class reachable? The oracle}
\label{sec:oracle}

Before attempting online adaptation we established that the fault is repairable
in principle by an edit at one site. A six-vector computed analytically was
added to \code{action\_out\_proj/bias}, scaled by $k$ (15 episodes per condition).

\begin{center}
\begin{tabular}{lrr}
\toprule
\textbf{Condition} & \textbf{Success} & \textbf{vs.\ floor} \\
\midrule
$k = -1.0$ (wrong sign), faulted            & 0.0\,\% (0/15)  & $-46.7$ \\
$k = 0.0$ (no edit), faulted \emph{(floor)} & 46.7\,\% (7/15) & --- \\
$k = 0.5$, faulted                          & 86.7\,\% (13/15) & $+40.0$ \\
\textbf{$k = 1.0$ (computed oracle), faulted} & \textbf{100.0\,\% (15/15)} & $\mathbf{+53.3}$ \\
$k = 1.5$, faulted                          & 20.0\,\% (3/15) & $-26.7$ \\
$k = 1.0$ on the \emph{healthy} policy      & 6.7\,\% (1/15)  & --- \\
\bottomrule
\end{tabular}
\end{center}

\begin{description}[leftmargin=1.2em,style=nextline,itemsep=3pt]
\item[Complete.] $46.7\% \rightarrow 100\%$ against a 99.0\% nominal --- the
entire deficit, from a fixed six-vector on one bias. No gradient, no fine-tuning.
\item[Specific, not a general improvement.] The same edit on the \emph{healthy}
policy collapses it from 99\% to 6.7\%; sign-flipped it annihilates the faulted
one (0\%). A compensating inverse, not a tonic.
\item[Narrow basin.] $k = 1.5$ scores 20\% --- \emph{worse than not repairing}.
Any search over this class must be scaled to $k \in [0.4, 1.2]$.
\end{description}

\paragraph{Why the edit must be computed per dimension.} ``Cancel a $+0.05$
offset with a $-0.05$ edit'' is wrong, and would have read as the class being
unreachable. $\pi_{0.5}$ emits \emph{normalised} actions;
\code{Unnormalize} runs afterwards, so with
$\mathrm{env} = \tfrac{\mathrm{norm}+1}{2}(q_{99}-q_{01}) + q_{01}$ each channel
carries its own scale, and those scales differ by \textbf{7$\times$}. A uniform
fault in environment space is wildly \textbf{anisotropic} in the model's units.
This same anisotropy turns out to decide the whole scope condition
(\S\ref{sec:scope}).

%==========================================================================
\section{Theory of the online law}

\subsection{The three objects}

An adaptive law needs an \textbf{error} to drive to zero, a \textbf{plant model}
predicting error-free behaviour, and the \textbf{map} from the unknown parameter
to that error. Get any wrong and the law diverges. All three failure modes were
observed here before being fixed (\S\ref{sec:failures}).

\subsection{Plant}

One action unit \emph{commands} 0.05\,m of translation. \textbf{It does not
achieve it}: regressing achieved on commanded displacement over fault-free
rollouts gives a static gain of \textbf{0.21--0.25}. The arm covers about a fifth
of the commanded delta inside one 50\,ms period, so using \code{output\_max} as
the plant is wrong by $5\times$ and an adaptive law built on it mistakes ordinary
controller lag for a fault. An FIR model over the last $K=6$ commands captures it:
\begin{equation}
y_t \;=\; \sum_{k=0}^{K} h_k\, a_{t-k} \;+\; c .
\end{equation}

\begin{center}
\begin{tabular}{lcccccc}
\toprule
$R^2$ & $dx$ & $dy$ & $dz$ & $d r_x$ & $d r_y$ & $d r_z$ \\
\midrule
static ($K=0$)       & 0.885 & 0.961 & 0.978 & 0.416 & 0.091 & 0.077 \\
FIR ($K=6$)          & 0.975 & 0.982 & 0.982 & 0.489 & 0.109 & 0.168 \\
\textbf{FIR, after SO(3) fix} & \textbf{0.975} & \textbf{0.982} & \textbf{0.982} & 0.522 & 0.336 & \textbf{0.972} \\
\bottomrule
\end{tabular}
\end{center}

\subsection{Error signal}

With $u_t = a_t + c_t + f$ the executed action (policy command, our correction,
unknown fault) and $\hat P$ the identified fault-free plant,
\begin{equation}
r_t \;=\; y_t - \hat P(a_t + c_t) \;\approx\; M f .
\end{equation}
$r_t$ is independent of $c_t$ \emph{when $\hat P = P$ exactly}, so the estimator
does not chase its own correction. \S\ref{sec:phantom} shows what happens when
that condition fails, and corrects this claim.

\subsection{The map}

$M = \partial(\text{achieved motion})/\partial f$, measured \textbf{open loop}: a
recorded command sequence replayed with and without a fault, so commands are
identical by construction and any motion difference is the fault propagating.
One axis at a time, $\pm 0.02$ central difference:

\begin{center}
\small
\begin{tabular}{lrrrrrr}
\toprule
 & \multicolumn{6}{c}{\textbf{fault applied to}} \\
\cmidrule(l){2-7}
\textbf{observed} & $dx$ & $dy$ & $dz$ & $d r_x$ & $d r_y$ & $d r_z$ \\
\midrule
$dx$    & $\mathbf{0.297}$ & 0.021 & $-0.112$ & $-0.029$ & 0.081 & $-0.006$ \\
$dy$    & 0.008 & $\mathbf{0.272}$ & 0.023 & $-0.034$ & 0.016 & 0.003 \\
$dz$    & 0.029 & 0.032 & $\mathbf{0.126}$ & $-0.037$ & 0.012 & $-0.003$ \\
$d r_x$ & $-0.003$ & $-0.001$ & $-0.001$ & $\mathbf{0.253}$ & $-0.004$ & 0.002 \\
$d r_y$ & 0.013 & 0.001 & $-0.001$ & $-0.007$ & $\mathbf{0.276}$ & $-0.000$ \\
$d r_z$ & 0.001 & 0.004 & $-0.003$ & 0.011 & 0.005 & $\mathbf{0.244}$ \\
\bottomrule
\end{tabular}
\end{center}

Diagonal 0.13--0.30 (the plant damps the fault 3--8$\times$, so the law needs
$\approx 4\times$ gain); off-diagonal mass 26\% of $|M|$, dominated by a real
$dx \!\leftrightarrow\! dz$ term; $\mathrm{cond}(M) = 3.0$.

\subsection{Orientation increments: a metric bug that corrupted four results}
\label{sec:so3}

Rotation change must be the \textbf{relative rotation}, not a difference of
axis--angle vectors:
$\omega = \mathrm{axis\_angle}(q_{t+1} \otimes \mathrm{conj}(q_t))$.
Axis--angle is a \emph{chart}, not a vector space: the same rotation has
representations differing by $2\pi$ and the chart is singular at $0$ and $\pi$.
Over 2000 random 0.05\,rad increments the relative rotation recovers truth to
$1.4\times 10^{-14}$; the difference of charts errs by up to \textbf{6.28} --- a
full wraparound.

\begin{tcolorbox}[colback=red!3,colframe=red!45!black,title=\textbf{Retraction (2026-08-27)}]
An earlier version reported the rotation block of $M$ as ``a swapped,
sign-flipped pair'' and argued a per-axis law would therefore diverge.
\textbf{That was wrong} --- an artifact of the chart difference. With the correct
metric the rotation axes have ordinary diagonal gains of 0.276 and 0.244,
off-diagonal mass falls from 0.59 to 0.26, and the divergence argument does not
apply. Withdrawn.
\end{tcolorbox}

The wrong metric corrupted \textbf{four separate results}: rotation appeared
unpredictable, $M$ appeared cross-coupled, the gain law's rotation estimates were
meaningless, and deliberate excitation made performance \emph{worse}. One bug.

\subsection{The law}

\begin{equation}
\fhat \leftarrow \mathrm{clip}\!\left(
  \fhat + \gamma\!\left( \frac{M^{-1} r_t}{1 + \|r_t\|^2/\rho^2} - \fhat \right),
  \pm 0.15 \right), \qquad c_t = -\fhat
\end{equation}
with a deadzone for $\|r_t\| < \delta$; $\gamma = 0.08$, $\rho = 0.15$,
$\delta = 0.008$. Each robustness term earned its place by fixing an observed
failure (\S\ref{sec:failures}), not by assumption.

\subsection{The fixed point of (4) is biased, and by how much}
\label{sec:bias}

Equation~(4) divides the \emph{estimate} rather than the \emph{step}. Because the FIR
history is built from $u = a_t + c_t$, the command we believe we sent, while the world
executes $a_t + c_t + f$, the residual is $r_t \approx M f$ --- independent of $\fhat$.
The normaliser is then a constant attenuation, and the fixed point of (4) is
%
\begin{equation}
\fhat^{\,\star} = \frac{f}{1 + \|Mf\|^2/\rho^2} \;\neq\; f .
\end{equation}
%
This is a real defect, and it is measurable rather than rhetorical. At the observed
$\|r\| = 0.034$ with $\rho = 0.15$ the attenuation is $1/1.051$, i.e.\ the estimate should
land \textbf{4.9\,\% low}. The separation test measures $+0.049$ against a true $0.050$,
\textbf{2\,\% low} --- same sign, same order. The bias is real and small, and it does
\emph{not} account for the gain-fault result, which \S\ref{sec:gain} attributes to
excitation instead.

The fix is to drive the innovation $e_t = r_t - M\fhat$ to zero and normalise the step:
%
\begin{equation}
\fhat \leftarrow \mathrm{clip}\!\left(
  \fhat + \gamma\,\frac{M^{-1} e_t}{1 + \|e_t\|^2/\rho^2}, \pm 0.15 \right),
  \qquad e_t = r_t - M\fhat .
\end{equation}
%
Its fixed point is $\fhat = M^{-1} r = f$ exactly. It also gives the estimate a way home:
if the fault clears, $e_t = -M\fhat$ is still non-zero, so $\fhat$ decays, whereas a
deadzone on raw $r_t$ can freeze a correction that is no longer needed. This is available
as \code{--law innov}; the legacy law remains the default so that every number in this
report stays reproducible, and the two are to be compared head-to-head rather than
swapped on the strength of the algebra alone.

%==========================================================================
\section{What the correction is actually doing}
\label{sec:phantom}

\subsection{\texorpdfstring{$\fhat$}{f-hat} alone cannot distinguish identification from bias}

Run the law with \textbf{no fault at all} and $\fhat$ does not go to zero. It
settles at $[+0.045, +0.014, +0.073, -0.009, +0.017, -0.001]$ --- on $dz$ that is
146\% of the magnitude of the real fault. The estimator reports a large offset
where there is nothing to correct.

The diagnostic that separates them is the \textbf{separation}: the estimate under
a fault minus the estimate with no fault, on matched episodes.

\begin{center}
\begin{tabular}{lrrrr}
\toprule
 & \textbf{no fault} & \textbf{fault 0.05} & \textbf{separation} & \textbf{true} \\
\midrule
$dx$    & $+0.045$ & $+0.048$ & $\mathbf{+0.003}$ & 0.050 \\
$dy$    & $+0.014$ & $+0.038$ & $+0.023$ & 0.050 \\
$dz$    & $+0.073$ & $+0.079$ & $\mathbf{+0.006}$ & 0.050 \\
$\mathbf{d r_x}$ & $-0.009$ & $+0.041$ & $\mathbf{+0.049}$ & 0.050 \\
$d r_y$ & $+0.017$ & $+0.033$ & $+0.016$ & 0.050 \\
$\mathbf{d r_z}$ & $-0.001$ & $+0.047$ & $\mathbf{+0.048}$ & 0.050 \\
\bottomrule
\end{tabular}
\end{center}

\textbf{Rotation identifies the fault almost exactly. Translation does not
identify it at all.}

\begin{tcolorbox}[colback=gray!6,colframe=gray!55!black]
\textbf{Method note.} No estimate of a fault parameter should be believed without
a matched no-fault control. $\fhat$ looked convincing on all six axes for weeks
of runs; only the separation revealed that half of it was bias.
\end{tcolorbox}

\subsection{Where the phantom comes from}

Measured separately with \code{--estimate-only} (update $\fhat$, never apply it):
mean $|$phantom$|$ is \textbf{0.0143} open loop and \textbf{0.0266} closed loop.

\emph{Half is plant-model error}, and it is not task-specific --- identifying the
plant on 10 tasks instead of 3 made the phantom \emph{worse} ($0.027 \to 0.034$).
\emph{Half is estimator feedback, which corrects the claim in \S5.3}: $r_t$ is
independent of $c_t$ only when $\hat P = P$ exactly. With model error
$r = Mf + \varepsilon(a+c)$, so $\fhat$ converges to $f + M^{-1}\varepsilon(a+c)$
--- it does chase its own correction. Measured amplification: \textbf{1.9$\times$}.

%==========================================================================
\section{The scope condition: correct only what you can identify}
\label{sec:scope}

\subsection{With all six channels corrected, the method works on one suite of four}

\begin{center}
\begin{tabular}{lrrrr}
\toprule
\textbf{Suite} & \textbf{frozen} & \textbf{corrected} & $\Delta$ & $p$ \\
\midrule
\code{libero\_spatial} & 18/40 = 45\% & 38/40 = 95\% & $+50$ & $\mathbf{<0.0001}$ \\
\code{libero\_goal}    & 8/20 = 40\%  & 12/20 = 60\% & $+20$ & 0.34 \\
\code{libero\_object}  & 7/20 = 35\%  & 9/20 = 45\%  & $+10$ & 0.75 \\
\code{libero\_10}      & 0/20 = 0\%   & 3/20 = 15\%  & $+15$ & 0.23 \\
\midrule
pooled, 3 new suites   & 15/60 = 25\% & 24/60 = 40\% & $+15$ & 0.12 \\
\bottomrule
\end{tabular}
\end{center}

It is \textbf{not} an identification failure: plant fits on the new suites are
comparable or better (\code{object} reaches rotation $R^2$ 0.771/0.692 against
spatial's 0.522/0.336). The estimator still works; the \emph{task} does not
recover.

\subsection{Restricted to identifiable channels, it works on all four}

\begin{center}
\begin{tabular}{lrrrr}
\toprule
\textbf{Suite} & \textbf{frozen} & \textbf{corrected} & $\Delta$ & $p$ \\
\midrule
\code{libero\_spatial} & 7/15 = 47\%  & 14/15 = 93\% & $+47$ & \textbf{0.014} \\
\code{libero\_goal}    & 9/20 = 45\%  & 18/20 = 90\% & $+45$ & \textbf{0.006} \\
\code{libero\_object}  & 6/20 = 30\%  & 16/20 = 80\% & $+50$ & \textbf{0.004} \\
\code{libero\_10}      & 0/20 = 0\%   & 7/20 = 35\%  & $+35$ & \textbf{0.008} \\
\midrule
\textbf{pooled} & \textbf{22/75 = 29\%} & \textbf{55/75 = 73\%} & $\mathbf{+44}$ & $\mathbf{1.1\times 10^{-7}}$ \\
\bottomrule
\end{tabular}
\end{center}

\code{libero\_10} is the sharpest case: a long-horizon suite (520-step cap) where
the fault takes the policy to \textbf{zero}, and correcting three channels
recovers 7/20 from that floor.

\subsection{Non-identifiable channels do not merely fail --- they do harm}

On \code{libero\_object}, three arms on the same fault and the same episodes:

\begin{center}
\begin{tabular}{lrrr}
\toprule
\textbf{Correction} & \textbf{frozen} & \textbf{corrected} & $\Delta$ \\
\midrule
all six dims          & 7/20 & 9/20  & $+10$ \\
\textbf{rotation only} & 6/20 & \textbf{16/20} & $\mathbf{+50}$ \\
\textbf{translation only} & 6/20 & \textbf{6/20} & $\mathbf{0}$ \\
\bottomrule
\end{tabular}
\end{center}

Translation alone contributes \textbf{exactly nothing}, and including it drags a
$+50$ effect down to $+10$. On \code{object} the translation estimates are
\textbf{sign-wrong} ($-0.031, -0.018$ against $+0.050$), so that correction pushes
the arm the wrong way. On \code{spatial} the same wrong correction happened to be
harmless --- which is why the naive version looked suite-specific rather than
simply misapplied.

\subsection{Identifiability is predictable before deployment}

Nothing here requires knowing the fault. Three quantities, all measurable on
\emph{healthy} data, decide which channels are identifiable:

\begin{enumerate}[leftmargin=1.5em,itemsep=2pt]
  \item \textbf{The action normalisation} --- quantile scales from the checkpoint.
        A uniform env-space offset is 3\% of the action range on translation and
        \textbf{19.5\%} on rotation, $6.6\times$ larger.
  \item \textbf{The plant fit} --- per-channel FIR $R^2$ on fault-free rollouts.
        Identification quality tracks it in every experiment run.
  \item \textbf{The policy's command statistics} --- an additive fault is loud
        where the quantile scale is small; a multiplicative one is loud where the
        command is large.
\end{enumerate}

\begin{tcolorbox}[colback=blue!4,colframe=blue!45!black]
\textbf{The rule.} Measure identifiability on healthy data, correct only those
channels, leave the rest alone. The all-channel table is the ablation showing
what it costs to ignore this.
\end{tcolorbox}

%==========================================================================
\section{Identifiability, stated properly}

\subsection{Proposition 1 --- a constant input fault and a constant output bias are not separable}

Suppose the measurement carries an unknown constant bias $b$:
$y_t = M(a_t + f) + b + \text{noise}$.
Within one episode $f$ and $b$ are both constant, so they enter only through the
sum $Mf + b$. $M$ is square and full rank, so for \emph{any} $b'$ the alternative
fault $f' = f + M^{-1}(b - b')$ reproduces every observation exactly. The pair
$(f, b)$ is unidentifiable; only $Mf + b$ is. \textbf{The bias must come from
outside the episode} --- it is not a nuisance parameter one can fit.

\paragraph{Corollary (onset breaks the degeneracy).} If the fault switches on at
$t_0$, with $f = 0$ for $t < t_0$, then $b$ is identified on the pre-onset window
and $f$ on the post-onset window. \textbf{Mid-episode onset is not a robustness
flourish --- it is what makes the problem well-posed}, and it is the honest
deployment story: the estimator needs to have seen the healthy plant.

\subsection{Proposition 2 --- the two fault classes are loud on complementary channels}

With quantile normalisation and $R_i = (q_{99}-q_{01})_i$: an \textbf{additive}
env-space offset $\delta$ appears in normalised units as $2\delta/R_i$, so its
signal \emph{falls} with $R_i$. A \textbf{multiplicative} fault has information
matrix $\sum_t \phi^\top\phi = \mathrm{diag}(\sum_t a_{t,i}^2)$, so its signal
\emph{rises} with command magnitude, and command spread is what $R_i$ measures.

\begin{center}
\begin{tabular}{lrrr}
\toprule
\textbf{ch} & $q_{99}-q_{01}$ & \textbf{additive signal} & \textbf{multiplicative signal} \\
\midrule
$x$  & 1.685 & 0.059 & 0.899 \\
$y$  & 1.656 & 0.060 & 0.883 \\
$z$  & 1.875 & 0.053 & 1.000 \\
$r_x$ & 0.256 & \textbf{0.391} & 0.137 \\
$r_y$ & 0.351 & \textbf{0.285} & 0.187 \\
$r_z$ & 0.506 & \textbf{0.198} & 0.270 \\
\midrule
translation mean & & 0.058 & \textbf{0.927} \\
\bottomrule
\end{tabular}
\end{center}

\textbf{The same quantity that makes a channel quiet for an offset makes it loud
for a gain.} This is a prediction, not a fit: it says an offset fault should be
correctable on rotation and a gain fault on translation, which is what
\S\ref{sec:scope} and \S\ref{sec:gain} find.

%==========================================================================
\section{Results}
\label{sec:results}

\subsection{Headline, confirmatory}

\begin{center}
\begin{tabular}{lrl}
\toprule
\textbf{Arm} & \textbf{Success} & \textbf{95\% CI} \\
\midrule
nominal (no fault) & 99\,\% & --- \\
frozen, faulted    & \textbf{18/40 = 45\,\%} & [31, 60] \\
\textbf{adaptive}  & \textbf{38/40 = 95\,\%} & [83, 99] \\
\bottomrule
\end{tabular}
\end{center}

$+50$ points, Fisher $p = 1.1\times 10^{-6}$, on initial states 40--43 that no
earlier run touched. The earlier $n{=}15$ estimate ($47\% \to 93\%$, $p = 0.014$)
did not shrink under retest --- it grew slightly, and the estimates reproduced
within $\pm 0.004$ on all six axes across two independent runs.

\paragraph{A caveat on every $p$-value in this report.}
The arms are \emph{paired}: both run the same $(\text{task}, \text{init})$ episodes with
the same policy and seeds, differing only in whether the correction is applied. Fisher's
exact test treats them as independent samples and therefore discards that pairing, so the
$p$-values quoted here answer a weaker question than the design supports. The correct
tests are exact McNemar on the discordant pairs, or a paired permutation test
(\code{openpi/mcnemar.py}). Conditioning on the pairs normally \emph{increases} power, so
these numbers are expected to be conservative rather than optimistic --- but that is an
expectation, not a result. The runs that produced the tables above recorded only per-arm
totals, not per-episode outcomes, so \textbf{the pairing cannot be recovered from the
stored files and the paired tests cannot be computed retrospectively}. Per-episode logging
is now in place; the headline conditions must be rerun before any paired $p$-value is
claimed.

\subsection{The full condition set}

\begin{center}
\begin{tabular}{lrrl}
\toprule
\textbf{Condition} & \textbf{frozen} & \textbf{adaptive} & \textbf{note} \\
\midrule
offset 0.05, from step 1 ($n{=}40$) & 45\% & \textbf{95\%} & $p = 1.1\text{e-}6$ \\
offset 0.10, near-lethal ($n{=}8$)  & 0\%  & 38\% & frozen fails every episode \\
\textbf{structured} $[0,0,0,{+}.06,{-}.06,{+}.03]$ ($n{=}12$) & 17\% & \textbf{83\%} & shape never tuned on \\
\textbf{onset at step 15} ($n{=}20$) & 60\% & \textbf{90\%} & $p = 0.065$; tracking \\
onset at step 45 ($n{=}15$) & 73\% & 80\% & little headroom by design \\
\textbf{no fault} ($n{=}15$) & 100\% & \textbf{100\%} & safe to leave running \\
gain 0.5, multiplicative ($n{=}10$) & 50--70\% & 70--90\% & noise-dominated \\
\bottomrule
\end{tabular}
\end{center}

\subsection{Generalisation to an untuned fault shape}

The structured fault $[0,0,0,+0.06,-0.06,+0.03]$ --- mixed signs, rotation only,
unseen magnitudes --- is recovered as a \emph{pattern}, not a scalar:

\begin{center}
\begin{tabular}{lrrr}
\toprule
\textbf{dim} & \textbf{true} & \textbf{separation} & \textbf{error} \\
\midrule
$d r_x$ & $+0.060$ & $+0.053$ & $-0.007$ \\
$d r_y$ & $-0.060$ & $-0.038$ & $+0.022$ \\
$d r_z$ & $+0.030$ & $+0.029$ & $-0.001$ \\
$dx, dy, dz$ & $0.000$ & $-0.003, -0.024, -0.006$ & $\approx 0$ \\
\bottomrule
\end{tabular}
\end{center}

Sign correct on 3/3 including the negative axis, and correctly $\approx 0$ where
the fault is zero. This rules out the objection that the method was fitted to one
fault shape.

\subsection{Tracking a fault that appears mid-episode}

The deployment case --- a healthy robot that breaks while running. Fault injected
at control step 15; mean $\fhat$ on identifiable rotation channels over 20
episodes:

\begin{center}
\begin{tabular}{lr}
\toprule
\textbf{window} & $\fhat$ \\
\midrule
steps 0--14, before onset & $\mathbf{+0.009}$ \\
steps 15--30, just after  & $+0.035$ \\
steps 30--50, settling    & $+0.041$ \\
steps 50--78, late        & $\mathbf{+0.042}$ (true 0.050) \\
\bottomrule
\end{tabular}
\end{center}

Quiet before the fault exists, \textbf{70\% of truth within 15 control steps
(0.75\,s)}, settling at 84\% and holding flat with no post-convergence drift.
This is also the cleanest control in the set: \textbf{the episode is its own
baseline}, so the estimator bias of \S\ref{sec:phantom} cancels without a
separate matched run.

%==========================================================================
\section{The gain fault: a state-dependent fault needs a regressor}
\label{sec:gain}

A constant offset is the easy case --- constant, so a constant correction cancels
it. A gain error $u = g\,a$ is state-dependent and needs a regressor: with
$\beta = g-1$, $z = M^{-1} r \approx \beta \odot \psi$, $\psi = a + c$.

Gain fault $g = 0.5$ (true $\beta = -0.50$), 10 episodes:

\begin{center}
\begin{tabular}{lccrrr}
\toprule
\textbf{Variant} & \textbf{Adaptive} & \textbf{Frozen} & \textbf{transl.\ $\beta$} & \textbf{rot.\ $\beta$} & \textbf{rot upd/ep} \\
\midrule
LMS, old metric            & 7/10 & 6/10 & $-0.589$ & $+0.016$ & 2 \\
RLS + dither, old metric   & 5/10 & 7/10 & $-0.617$ & $+0.051$ & 136 \\
\textbf{RLS + dither, SO(3) fix} & \textbf{7/10} & 5/10 & $-0.636$ & $-0.169$ & 30 \\
\bottomrule
\end{tabular}
\end{center}

Per-dimension $\beta$ after the fix: $dx\ {-}0.62$, $dy\ {-}0.49$, $dz\ {-}0.80$,
$d r_x\ {-}0.20$, $d r_y\ {+}0.08$, $d r_z\ {-}0.38$. The metric fix does for the
gain law what it did for the offset law: $d r_z$ flips from $+0.363$ ---
confidently the \emph{wrong sign} --- to $-0.384$; $dy$ lands at $-0.49$ against
$-0.50$. Stated precisely, because the looser phrasing conflated two
different criteria: \textbf{five of six carry the right sign} ($d r_y$ at $+0.08$ is
the exception), but \textbf{only one of six, $dy$, is within $\pm 20\,\%$ of the true
$-0.50$}. The ratios $\hat\beta/\beta^\star$ are $1.24$, $0.98$, $1.60$, $0.40$,
$-0.16$, $0.76$. Sign recovery is not magnitude recovery, and only the sign claim
survives at this sample size.

\begin{tcolorbox}[colback=orange!5,colframe=orange!55!black,title=\textbf{Honest reading}]
Frozen scored 5, 6, and 7 of 10 across these three runs. \textbf{None of the
adaptive-vs-frozen success differences at $n{=}10$ are resolvable.} The $\beta$
columns carry the signal; the success columns are noise at this sample size.
\end{tcolorbox}

Consistent with Proposition 2, the gain fault is identified best on
\emph{translation} --- the complement of the offset fault's rotation channels.

%==========================================================================
\section{What this fault is, and what it is not}
\label{sec:whatfault}

The injected fault adds a constant to the six Cartesian action dimensions the policy
emits, upstream of the OSC controller. It is therefore a \textbf{Cartesian
action-interface bias}. It is a faithful model of anything that corrupts the commanded
end-effector increment --- a miscalibrated tool frame, a wrist-mount offset, a stale
hand--eye calibration --- and the recovery result is genuine for that class.

It is \emph{not} a demonstrated recovery from an actuator or joint-level hardware fault.
Such a fault would enter \emph{below} the OSC, in joint dynamics, where the controller
partially absorbs it and the map from fault to observed motion is state-dependent rather
than constant. The hardware work in this report shows exactly that gap: a joint stiffness
fault is invisible in position and recoverable only in effort. Claiming actuator-fault
recovery requires injecting the fault in joint space and rerunning; until then the honest
scope is the action interface, and language throughout has been changed accordingly.

A second scoping point. The correction is applied by editing \code{action\_out\_proj/bias}
inside the network. Because $\pi_{0.5}$ generates actions by flow-matching integration,
that edit is \emph{not} identically an addition of $c_t$ to the final action; the mapping
runs through the integrator and is only approximately affine. The report's estimates are
consistent with a near-affine mapping, but this has not been measured directly, and an
external adapter applied after unnormalisation would make the mapping exact and
architecture-independent. That is the recommended construction going forward.

%==========================================================================
\section{Negative results worth keeping}
\label{sec:failures}

\paragraph{Naive error signal, no plant model.} Using
$\Delta x/\code{output\_max} - a$ gives a nominal error of $-0.175$ on $dx$ when
it should be $0$. Differencing faulted against nominal recovered ratios of
$0.39$, $2.46$, $-2.85$ --- noise. \emph{Controller lag reads as a fault unless
the plant is modelled.}

\paragraph{Robustness constants set without measuring.} The bare law drifted ---
$dz$ reached $0.642$, \textbf{13$\times$ truth}. The first fix made it
\emph{worse}: a deadzone of $0.05$ against a typical residual norm of $0.034$
suppressed nearly every update, ending several episodes at $\fhat = 0$ exactly.
Rescaling to the measured residual fixed both; raising the normalisation constant
($0.05 \to 0.15$, cutting a systematic 32\% attenuation to $\approx 5\%$) took the
result from 70\% to 87\%.

\paragraph{A coupled MIMO plant: better fit, worse extrapolation.} Wrist rotation
is driven substantially by \emph{translation} commands, which a per-axis model
cannot represent. A coupled plant improves leave-one-episode-out held-out $R^2$
from $0.107 \to 0.615$ on $d r_y$. \textbf{And it made the closed loop worse}:
13/15 ($p = 0.109$) with $dy$ estimated at $-0.043$ against a true $+0.050$ --- a
sign error the per-axis model never makes --- versus 14/15 ($p = 0.014$). Forty-three
parameters fitted on three nominal episodes extrapolate badly once the correction
moves the executed command off that distribution, and \emph{the adaptive law lives
out of distribution by construction}. Kept as \code{--mimo}, default off.

\paragraph{Excitation is not the limiting factor for rotation.} The policy
essentially never commands rotation: translation command sd is 0.32--0.55 with
49--74\% of steps above 0.15, while all three rotation axes sit at sd 0.02--0.055
with \textbf{0.0\%} above 0.15. But excitation alone does not explain $d r_y$: it
has the most rotation excitation and the worst fit, while $d r_z$ has less and
fits at 0.972. Coupling, not excitation, is the cause.

\begin{tcolorbox}[colback=gray!6,colframe=gray!55!black]
\textbf{The recurring lesson.} An improvement measured in one regime --- open-loop
fit, or excitation --- can cost you in the regime that actually matters. The MIMO
plant and the dither result have the same shape.
\end{tcolorbox}

%==========================================================================
\section{The video}

\code{results/phase05/adaptive\_vs\_frozen.mp4} shows two runs side by side: the
same frozen model, the same task, the same injected offset fault. The only
difference is whether the online repair is active. One run is the frozen policy
--- the fault has shifted its actions and it fails the task. The other is the same
weights with the adaptive law running: it identifies the fault from the robot's
own motion within the first steps and cancels it, and the task completes normally.
Nothing is retrained between the two runs.

A companion clip, \code{adaptive\_vs\_frozen\_offset010.mp4}, shows the harder
severity (offset 0.10), where frozen success is 0/8.

%==========================================================================
\section{Limits and what this does not establish}
\label{sec:limits}

\begin{itemize}[leftmargin=1.4em,itemsep=3pt]
  \item \textbf{Simulation.} Four LIBERO suites; hardware validation is in
        progress and has already produced one negative --- the hardware plant is
        run-varying on the load-bearing joints, so the offline identification
        recipe will not port unchanged.
  \item \textbf{The law is blind to sensor bias.} $\pi_{0.5}$ ignores
        proprioception entirely, so a proprioceptive fault produces no change in
        the policy's behaviour for the estimator to detect.
  \item \textbf{Gain faults are not settled.} One severity, $n{=}10$, no
        resolvable success difference; only the $\beta$ estimates carry signal.
  \item \textbf{Rotation remains partly unidentified.} $d r_y$ still carries the
        wrong sign in the gain law and fits at $R^2 = 0.336$.
  \item \textbf{$M$ is identified open loop} and assumed valid in closed loop,
        where the policy compensates and the arm visits different states.
  \item \textbf{A gap to the oracle remains.} The oracle reaches 100\%, the
        adaptive law 95\% on \code{spatial} and 73\% pooled.
  \item \textbf{The site screen does not separate sites} at the sample size run;
        \code{action\_out\_proj/bias} was chosen structurally.
\end{itemize}

%==========================================================================
\section{Summary}

\begin{enumerate}[leftmargin=1.6em,itemsep=3pt]
  \item A frozen 3.35\,B VLA under a miscalibrated-actuator fault drops from 99\%
        to 45\% task success.
  \item The fault is \textbf{exactly repairable} by a six-number edit to one
        layer's bias ($47\% \to 100\%$), provided the edit is computed
        \emph{per dimension in the model's normalised coordinates}, where a
        uniform physical fault is anisotropic by $7\times$.
  \item The repair can be \textbf{identified online from proprioception alone},
        within a single episode, with no gradients and no search over task
        success --- and it tracks a fault that appears mid-episode to 70\% of
        truth in 0.75\,s.
  \item \textbf{The contribution is a scope condition.} Across four suites the
        method is significant on one of four when all channels are corrected, and
        on \emph{all four} ($29\% \to 73\%$, $p = 1.1\times10^{-7}$) when
        restricted to identifiable channels. Correcting a non-identifiable
        channel is \emph{worse than not correcting at all}.
  \item \textbf{Identifiability is decidable in advance} from healthy data ---
        action normalisation, plant fit, and command statistics --- and the theory
        predicts that additive and multiplicative faults are identifiable on
        \emph{complementary} channels, which the experiments bear out.
  \item Getting there required a metric bug on SO(3) that had silently corrupted
        four results, and a matched no-fault control that revealed half of a
        convincing-looking estimate was bias.
\end{enumerate}

\end{document}
